\subsubsection{Exact logarithmic-square growth of the central heat source}
\label{subsec:theta-central-asymptotic}
The divergence established in
Theorem~\ref{thm:centralmomentratio}
has a definite leading order.
This can be obtained without knowing the zeros of
$H_t$ and without applying a saddle-point formula:
the decisive input is the double-exponential
tail of the \emph{actual} theta convolution.

\begin{theorem}[Uniform central moment-ratio asymptotic]
\label{thm:theta-central-logsquare}
Let $\Phi$ be the even classical Jacobi theta kernel
of this paper, and set
\[
\Phi_t(u)=e^{tu^2}\Phi(u),\quad
g_t(w)=(\Phi_t*\Phi_t)(w),\quad
M_{2n}(t)=\int_{\mathbb R}w^{2n}g_t(w)\,dw
\]
for $0\le t\le1/2$. Then
\begin{equation}\label{eq:uniformmomentasymptotic}
\boxed{\lim_{n\to\infty}\ \sup_{0\le t\le1/2}
\left|
\frac{M_{2n+2}(t)}
{M_{2n}(t)(\log n)^2}-\frac14\right|=0.}
\end{equation}
Consequently, with $U_n(s,x)=L_{n,1/2-s}(x)$
and $k_n=(n+1)(2n+1)$,
\begin{equation}\label{eq:centralexactgrowth}
\boxed{
\frac{k_nU_{n+1}(s,0)}{U_n(s,0)}
=\left(\frac18+o(1)\right)(\log n)^2}
\end{equation}
uniformly for $0\le s\le1/2$.
The $o(1)$ is uniform in heat time.
\end{theorem}

\begin{proof}
We first derive uniform two-sided tail estimates.
For $u\ge0$, the single $m=1$ summand and $t\ge0$
give, with $q_1=\pi$,
\[
\Phi_t(u)\ge
\pi(2\pi-3)e^{-\pi e^{4u}}.
\]
Conversely, polynomial factors can be absorbed into
a fixed fraction of the first exponential.
The remaining sum over square-lattice indices
is uniformly convergent for $e^{4u}\ge1$.
Since $0\le t\le1/2$ and $e^{4u}$
dominates $u^2$, there are positive constants
$a,A,c,C$, independent of $t$, such that
\begin{equation}\label{eq:theta-doubleexp}
a e^{-A e^{4|u|}}
\le\Phi_t(u)\le C e^{-c e^{4|u|}}
\qquad (u\in\mathbb R,\ 0\le t\le1/2).
\end{equation}
Here the bounds for negative $u$ use the exact
Jacobi evenness of $\Phi$.

For $w\ge1$, restrict the convolution integral
to $u\in[w/2-1/4,w/2+1/4]$. Both $u$ and $w-u$
then have absolute value at most $w/2+1/4$.
The lower bound in \eqref{eq:theta-doubleexp}
therefore yields
\[
g_t(w)\ge
\frac{a^2}{2}
\exp(-2Ae\,e^{2w}).
\]
For the upper bound, at least one of $|u|$ and
$|w-u|$ is not less than $w/2$.
The other factor is integrable, with its $L^1$
norm bounded uniformly in $t$.
This gives constants $a_0,A_0,c_0,C_0>0$,
independent of heat time, with
\begin{equation}\label{eq:convolution-doubleexp}
a_0e^{-A_0 e^{2w}}
\le g_t(w)\le C_0e^{-c_0 e^{2w}},
\qquad w\ge1.
\end{equation}
The density is even, positive, and has finite
moments of every order.

Write $I_n(t)=\int_0^\infty w^{2n}g_t(w)\,dw$.
Fix $0<\varepsilon<1$, and put
\[
\alpha=\frac{1-\varepsilon}{2},\qquad
\beta=\frac{1+\varepsilon}{2},\qquad
\gamma=1-\frac{\varepsilon}{2},\qquad
v_n=\frac\gamma2\log n.
\]
Thus $\alpha<\gamma/2<\beta$.
Integrating the lower bound of
\eqref{eq:convolution-doubleexp} on
$[v_n,v_n+1]$ gives, uniformly in $t$,
\begin{equation}\label{eq:denom-lower}
I_n(t)\ge
a_0 v_n^{2n}\exp(-A_0e^2 n^\gamma).
\end{equation}

Set $a_n=\alpha\log n$ and $b_n=\beta\log n$.
Since $g_t$ is uniformly bounded, the
small-$w$ portion of the next moment obeys
\[
\int_0^{a_n}w^{2n+2}g_t(w)\,dw
\le C_1 a_n^{2n+3}.
\]
Relative to $(\log n)^2 I_n(t)$,
this is bounded by a polynomial in $\log n$
times
\[
\left(\frac{2\alpha}{\gamma}\right)^{2n}
\exp(A_0e^2n^\gamma),
\]
which tends to zero because $2\alpha<\gamma<1$.

For $w\ge b_n$ and all sufficiently large $n$,
the inequality
\[
(2n+2)\log w\le \frac{c_0}{2}e^{2w}
\]
holds throughout the half-line: it holds at
$b_n$, since $e^{2b_n}=n^{1+\varepsilon}$,
and the difference is increasing thereafter.
By \eqref{eq:convolution-doubleexp},
\[
\int_{b_n}^\infty w^{2n+2}g_t(w)\,dw
\le C_0\int_{b_n}^\infty
e^{-(c_0/2)e^{2w}}\,dw
\le C_2 e^{-(c_0/4)n^{1+\varepsilon}}.
\]
This is $o((\log n)^2 I_n(t))$
by \eqref{eq:denom-lower}, again uniformly
in $t$. The same estimates, with the power
$2n$ instead of $2n+2$, show that the
$w^{2n}g_t(w)$-probability of the complement
of $[a_n,b_n]$ tends to zero uniformly.

It follows that
\[
\alpha^2-o(1)\le
\frac{I_{n+1}(t)}{I_n(t)(\log n)^2}
\le\beta^2+o(1)
\]
uniformly for $0\le t\le1/2$.
Letting $\varepsilon\downarrow0$ proves
\eqref{eq:uniformmomentasymptotic};
the factor two from integrating over
positive and negative $w$ cancels in
the moment ratio.

Finally, the exact central identity
\eqref{eq:centralratio} gives
$k_nU_{n+1}(s,0)/U_n(s,0)
=\frac12M_{2n+2}(t)/M_{2n}(t)$,
where $t=1/2-s$, proving
\eqref{eq:centralexactgrowth}.
\end{proof}

\begin{corollary}[Necessary order dependence of a
global spatial heat comparison]\label{cor:necessary-order}
Suppose nonnegative coefficients $c_n(s,x)$
satisfy
\[
k_nU_{n+1}(s,x)\le c_n(s,x)U_n(s,x)
\quad
(n\ge1,\ 0\le s\le1/2,\ x\in\mathbb R).
\]
Then necessarily
\begin{equation}\label{eq:necessarycn}
\liminf_{n\to\infty}\ 
\inf_{0\le s\le1/2}
\frac{c_n(s,0)}{(\log n)^2}\ge\frac18.
\end{equation}
In particular no coefficient with
$c_n(s,0)=o((\log n)^2)$ uniformly
in $s$ can satisfy the required comparison.
\end{corollary}
\begin{proof}
The positive central convolution
gives $U_n(s,0)>0$ for every $n,s$.
Divide the proposed source inequality
by $U_n(s,0)$ and invoke
\eqref{eq:centralexactgrowth}.
\end{proof}

The scale $(\log n)^2$ is an \emph{order}
requirement at the centre, distinct from
the $(\log|x|)^2$ lower-growth
obstruction arising from increasing zero
density at positive heat time.
The simultaneous spatial-and-order
upper comparison remains unproved;
a prospective envelope such as
$C\{(\log(n+2))^2+(\log(e+|x|))^2\}$
must not be asserted without controlling
zero collisions and the noncentral
oscillatory correlation.
Moreover, the asymptotic alone is not a
sufficient criterion for RH.

The two-sided tail argument above is analytic,
not an end-to-end Lean formalization.
Exact finite algebraic transfer gates are in
\path{XiKernel/ThetaCentralLogSquare.lean};
the accompanying tests check these gates,
not the infinite asymptotic statement.
The novelty of the asymptotic relative to
earlier moment-asymptotic literature has
not yet been independently established.
